% Pursuing Stacks (cotes 134-2 to 134-8): archive cote 134-6, site page 86 (PDF page 87), typescript page 404.
% typescript. Draft compared: ps5.tex:5849-5911.
% AI-transcribed by Claude Opus 5.5 (claude-opus-5-5, adaptive thinking, effort high), Message Batches API, batch ps3 (msgbatch_01XdFmuVVBAcgc7EFxdtCqsj), 2026-11-12; not yet checked by eye.
\scanpage{134-6/86}
\mspage{404}

$P^{\&}$, $Q^{\&}$ and the category of abelian cosheaves on the other.

We may call an abelian topos ``\uline{reflexive}'' if it can be inserted in a
pair ($\hand{\mathcal{P}}$,$\hand{\mathcal{Q}}$) of dually paired \xout{topoi} abelian topoi -- where $\hand{\mathcal{Q}}$, or the
``dual'' of $\hand{\mathcal{P}}$, is defined up to equivalence in terms of $\hand{\mathcal{P}}$ as $\underline{\text{Hom}}_{\hand{!}}(\hand{\mathcal{P}},(\text{Ab}))$\hand{,}
the category of cosheaves on $\hand{\mathcal{P}}$ with values in (Ab). Thus, it seems that
a sufficient condition \xout{x}for reflexivity is ``algebraicity'' of $\hand{\mathcal{P}}$, namely
the existence of a small generating family \xout{of} made up with ultrapro\hand{j}ective
objects. (NB In the non abelian case, it is well known that a topos
$\hand{\mathcal{A}}$ is ``algebraic'' i.e. equivalent to a topos $A^{\wedge}$\hand{,} iff it admits such a
generating family -- and as we saw in 1º), as a consequence of (39)
for $\hand{\mathcal{M}}$= (Sets), such a topos is indeed ``reflexive'' .) I would'nt be
too surprised if this \xout{necessary}\inserted[typed]{sufficient}condition for reflexivity turned out to
be necessary too, at any rate if we want a property stronger still than
reflexivity, namely validity of a duality equivalence (\hand{***)}for sheaves
and cosheaves with values in an arbitrary additive $\underline{U}$-category stable under
small direct and inverse limits, satisfying (for varying $\hand{\mathcal{M}}$) suitable
compatibility assumptions.

\hand{3)}With respect to this duality equivalence (\hand{***)}\hand{,} I am a little unhappy
still, as I do not see how to get (for a general dual pair $\hand{\mathcal{P}}$, $\hand{\mathcal{Q}}$ of
abelian topoi) a functor in one direction or the other between the two
categories
\[ \underline{\text{Hom}}_{\hand{!}}(\hand{\mathcal{P}},\hand{\mathcal{M}}) \quad , \quad \underline{\text{Hom}}^{\hand{!}}(\hand{\mathcal{Q}^{o}},\hand{\mathcal{M}}) \quad , \]
in terms of just the duality pairing. The same perplexity holds in the
non abelian case. This is one of the reasons which make me feel that I
hav'nt yet a thorough understanding of the duality formalism I am develop-%
ping here, except in the ``algebraic'' case (granting \inserted{for the latter,} that the tentative
theory just outlined for algebraic \uline{abelian} topoi is indeed correct).

\hand{4)} To finish with these comments on the (pre-homological) duality for-%
malism for algebraic topoi and algebraic abelian topoi, I still would
like to add that the category $\hand{\mathcal{E}}$(A) (union of $A^{\wedge}$ and $B^{\wedge o}$, \xout{xxx}\inserted[typed]{with} $B=A^{o}$ )
introduced in 1º) (cf. remark on page \hand{\unsure{398}}) admits also an abelian
analogon. In the non abelian case still, the simplest way to construct
the category $\hand{\mathcal{E}}$(A ), is via an equivalent category canonically embedded
in the category $\underline{\text{Hom}}(B^{\wedge}, (\text{Sets}))$ as a strictly full subcategory
(``strict'' referring to the fact that \xout{xxxobject} with any object it contai\hand{ns}
all isomorphic ones), namely the union $\hand{\overline{\mathcal{E}}(A)}$ of the \inserted{(strictly full)} subcategories

\margin{86}
